Input completeness and output meaning must also be explicit. A model that
fixes the available event templates and contract slots imposes a scope
restriction; it does not establish that a public filing lists every economic
event. Nor does that restriction entitle a display to interpret every omitted
row as false.

Libkin's primary proof compares information by reverse inclusion of possible
completions: an object is more informative when fewer complete objects remain
compatible with it. For a query to be evaluated directly on an incomplete
object, the chosen output interpretation must preserve that order as input
information increases. The query must also respect the relevant renamings,
and its output domain must have a representation system whose formulas
describe each object's complete interpretations
\citep[Sections~6.1--6.2, equations~(9)--(10)]{libkin2014}. In Theorem~2, the proof
first shows that evaluating the query preserves exactly the statements true
in all completed answers; it then uses those statements to establish the
greatest-lower-bound answer as an object
\citep[Definitions~4--7 and Theorem~2, pp.~11--18]{libkin2016certain}.
Monotonicity alone is insufficient, and SQL null comparisons need not meet
the genericity condition \citep[Section~7.1, pp.~18--19]{libkin2016certain}.

The relational examples make the difference between input and output
assumptions concrete. Positive existential queries have the stated result
with open-world inputs and answers. A guarded positive fragment also has it
with \emph{closed-world inputs and open-world answers}; the latter is an
interpretation of omitted result rows, not a claim that the input filing
lists every economic event
\citep[Proposition~9 and Corollary~3, p.~20]{libkin2016certain}.
These results do not authorize evaluating arbitrary financial negation,
aggregates or nonlinear constraints by treating unknown amounts as ordinary
constants. AccDSL has not established the theorem's abstract domain,
information order and representation-system conditions for every financial
output. The constructed payable example above instead follows directly
from its two declared histories.
